\color{unidarkgreen}{\subsection{Логика опробования}}
\color{black}

\begin{enumerate}[label=\arabic{section}.\arabic{subsection}.\arabic{enumi}, labelsep=4pt, leftmargin=0pt, itemindent=57pt, itemsep=0pt, parsep=5pt] % если что можно крутить itemsep

%%%%%%%%%%%%%%%%%%%%%%%%%%%%%%%%%%%%%%%%%%%%%%%%%% Логика опробования %%%%%%%%%%%%%%%%%%%%%%%%%%%%%%%%%%%%%%%%%%%%%%%%%%
\item Общие данные

\begin{enumerate}[label=\arabic{section}.\arabic{subsection}.\arabic{enumi}.\arabic*, labelsep=4pt, leftmargin=0em, itemindent=65pt, parsep=0pt]

\item
Логика опробования предназначена для работы ДЗШ при оперативном опробовании выключателей присоединений по схеме с «открытым» или «закрытым» плечом.
При неуспешном опробовании отключается только опробуемый выключатель, остальные выключатели остаются в работе.

\medskip

\begin{figure}[H]
\centering
\includegraphics[width=0.83\textwidth,height=0.83\textheight,keepaspectratio]{\fbpath/661.1371 Логика опробования/img/Log_oprob.pdf}
\captionof{figure}{Функциональный блок логики опробования}\label{pic:Log_oprob}
\end{figure}

\medskip

\item
Переход в режим опробования осуществляется при наличии разрешающего сигнала \textit{«Ввод опробования»}.
Запуск опробования выполняется при поступлении команды РКВ от выключателя любого присоединения, зафиксированного за определенной СШ (при наличии сигнала \textit{«Фиксация Вn 1С»} или \textit{«Фиксация Вn 2С»}, где n – номер присоединения).
Продолжительность опробования определяется уставкой выдержки времени \textbf{«Т опроб»}.

\item
С помощью программного переключателя \textbf{«Опроб. Вn»} задается условие для отключения опробуемого выключателя по срабатыванию одного из следующих органов: ДЗШ ПО, ЧТО ПО или индивидуального ИО по току опробуемого присоединения с блокировкой по второй гармонике.

\item
Ток срабатывания индивидуального ИО задается с помощью уставки \textbf{«I опроб. Вn»}.
Блокировка по второй гармонике необходима для предотвращения ложных срабатываний ИО при БНТ трансформатора.
Для ввода блокировки предусмотрен программный ключ \textbf{«БНТ Вn»}. Уровень срабатывания блокировки задается с помощью уставки \textbf{«I2г/ I1г Вn»}.
Срабатывание ИО осуществляется по завершении выдержки времени \textbf{«Тср. опроб. Bn»}.

\item
Для обеспечения отключения только опробуемого выключателя, без отключения всей СШ, предусмотрены сигналы блокировки ДЗШ: \textit{«Блок. ДЗШ 1С при опроб.»} (\textit{«Блок. ДЗШ 2С при опроб.»} и \textit{«Блок. ДЗШ при опроб.»}).
Время действия блокировки соответствует продолжительности опробования.

\item
При нарушении фиксации присоединений сигналы блокировки ДЗШ формируются для обеих СШ.

\item
Уставки функции опробования представлены в таблице \ref{desc:tbl1}.

\small
\begin{longtable}{|>{\centering\arraybackslash}m{5.3cm}|>{\centering\arraybackslash}m{3.3cm}|>{\centering\arraybackslash}m{4.2cm}|>{\centering\arraybackslash}m{1.8cm}|>{\centering\arraybackslash}m{1cm}|}
% Первая страница: полный заголовок
\caption{Перечень уставок функции опробования\hfill\vspace{-0.5\baselineskip}}\label{desc:tbl1}\\
\hline
\rowcolor{gray!30}
Параметр & Условное обозначение на схеме & Значение/ Диапазон & Единица измерения & Шаг \\
\hline
\endfirsthead

% Промежуточные страницы: «Продолжение таблицы»
\caption*{\hspace{3pt}\emph{Продолжение таблицы \ref{desc:tbl1}}\hfill\ }\vspace{-0.5\baselineskip} \\
\hline
\rowcolor{gray!30}
Параметр & Условное обозначение на схеме & Значение/ Диапазон & Единица измерения & Шаг \\
\hline
\endhead

%===m>ДЗШ|-
%===m
\end{longtable}
\normalsize	

\end{enumerate}

\end{enumerate}
